\documentclass{article}
\usepackage{amsmath}
\usepackage{cancel}
\usepackage{graphicx} % for creating figures

\newcommand{\prob}{\mathcal{P}}


\title{Temperature}

\begin{document}
\maketitle

We have established that at thermal equilibrium, two systems than can exchange
energy must satisfy
\begin{equation}
  \frac{\partial S_1}{\partial U_1} = \frac{\partial S_2}{\partial U_2}
\end{equation}
where
\begin{equation}
  S \equiv k_b\ln g(N,U).
\end{equation}
This directly leads to the idea of temperature. Our intuitive notion of
temperature is \textit{the thing that's the same when two objects are in thermal
  equilibrium}. In other words,
\begin{equation}
  T_1 = T_2 \qquad \textrm{(at equilibrium)}
\end{equation}
We also know that if I dump energy into the system (ex: microwave a bowl a
soup), then the temperature should increase. Putting this all together, we
conclude that a reasonable definition of temperature is
\begin{equation}
  \frac{1}{T} = \left( \frac{\partial S}{\partial U}\right)_N
\end{equation}






\end{document}
